\section*{Chapter \ref{ch:in:education-pipelines} --- Education Pipelines}

\begin{solution}{exo:in:education-pipelines:1}
$4\,584/305=15.0$: fifteen times.
\end{solution}

\begin{solution}{exo:in:education-pipelines:2}
$67/70=95.7\%$.
\end{solution}

\begin{solution}{exo:in:education-pipelines:3}
42: six problems of 7 points each.
\end{solution}

\begin{solution}{exo:in:education-pipelines:4}
With 38 non-reporters, between the $(49-38)/60=18.3$rd and $49/60=81.7$th percentiles of the reported pay: the
reported median would say little about the cohort's.
\end{solution}

\begin{solution}{exo:in:education-pipelines:5}
Programmes choose their code, and some financial engineering programmes use others (engineering, business, statistics),
so the code misses them; if programmes moved into the code over time, the count grows faster than the number of such
degrees.
\end{solution}

\begin{solution}{exo:in:education-pipelines:6}
$1\,760+1\,628+1\,165+454=5\,007$ research doctorates in physics, computer science, mathematics and statistics, against
$4\,584+870=5\,454$ master's degrees in the two financial codes: about as many.
\end{solution}

\begin{solution}{exo:in:education-pipelines:7}
Computer science master's degrees rose from 11\,598 to 25\,543, 2.20 times; financial mathematics from 3\,554 to 4\,584,
1.29 times.
\end{solution}

\begin{solution}{exo:in:education-pipelines:8}
The two reports use different denominators and standards: one does not say how many reported pay, the other counts
short-term employment as placement under its standard; the cohorts graduate at different times of year and into
different markets; and a median of base pay says nothing about bonus or the jobs' kinds. A \$5\,000 difference is within
what those choices can move.
\end{solution}

\begin{solution}{pb:in:education-pipelines:1}
\begin{enumerate}
\item As in the chapter's definition.
\item 27.0305 financial mathematics (``the application of mathematics and statistics to the finance industry''), and
30.7104 financial analytics (``financial big data modeling''), new in 2020.
\item Awards by institution, code, major and level, July to June.
\item So that a student with two majors is counted once.
\item Other programmes use other codes.
\item 305, 773, 3\,554 and 4\,584.
\item As in the degrees box.
\item Less than a fifth of computer science's master's degrees.
\item 1\,760, 1\,628, 1\,165 and 454: 5\,007.
\item Research from doctorates and general master's; engineering from computer science; bank quant and structuring also
from the specialised master's.
\item As in the chapter's definition.
\item 70 seeking, 67 accepted, \$150\,000, reporting count not stated; 98 seeking, 98 accepted within three months
(including short-term), 89 reporting, \$145\,000.
\item Between the reported quantiles $(n/2-(n-r))/r$ and $(n/2)/r$: 44.9\% and 55.1\% for 89 of 98.
\item What non-reporters earned and why they did not report; bonuses; what counts as a job.
\item Speed and depth on well-posed problems; not research on ill-posed ones.
\item 5\,454 master's degrees in 2024 against 1\,466 filings in fiscal 2025: 3.72.
\item The two counts cover different people: all graduates of all nationalities against visa filings by a sample of
employers.
\item Graduates of computer science, statistics, mathematics and physics as well as the other financial programmes.
\item Its denominators (class size, seeking, accepted, reporting), its reporting standard and the dates.
\item The market is several thousand specialised master's graduates a year beside tens of thousands of general ones, and
the firms hire from both. A new programme will be judged by placement evidence it can show with its denominators.
\end{enumerate}
\end{solution}

\begin{solution}{iq:in:education-pipelines:1}
A specific example of working on a problem without a known answer for months: dead ends, checking one's own results,
and deciding what to give up.
\iqlookfor{evidence of independent research.}
\end{solution}

\begin{solution}{iq:in:education-pipelines:2}
Present a project that can be rerun, and explain a design choice, what was tried, and what was measured to decide.
\iqlookfor{reasoning about trade-offs, not a list of tools.}
\end{solution}

\begin{solution}{iq:in:education-pipelines:3}
Of how many, how many sought work, what counts as a placement (short-term work?), by which date, and how many reported
pay.
\iqlookfor{denominators and definitions.}
\end{solution}

\begin{solution}{iq:in:education-pipelines:4}
A Fermi estimate from the supply side (graduates in feeding fields times the share hired) and from the demand side
(firms' headcounts times turnover and growth), with the two compared; state the uncertainty.
\iqlookfor{two independent estimates.}
\end{solution}

\begin{solution}{iq:in:education-pipelines:5}
With code others can read and run: a repository with tests, a build, documentation and history, and one piece you can
explain line by line.
\iqlookfor{evidence over assertion.}
\end{solution}

\begin{solution}{iq:in:education-pipelines:6}
Look for a matching fall in the codes it replaced, compare institutions that reported under both, and check totals that
do not depend on the code.
\iqlookfor{crosswalks and totals.}
\end{solution}
